\chapter{Internship: Pengujian Unit dan Modul Proyek}

\section{Identitas Mata Kuliah}
\begingroup
\renewcommand{\arraystretch}{1.15}
\noindent\begin{tabular}{@{}p{3.6cm}@{\hspace{0.2cm}}p{0.2cm}@{\hspace{0.2cm}}p{7.6cm}@{}}
    \textbf{Kode Mata Kuliah} & : & MII21-3015 \\
    \textbf{Nama Mata Kuliah} & : & Internship: Pengujian Unit dan Modul Proyek \\
    \textbf{Bobot} & : & 3 SKS \\
    \textbf{Bentuk Kegiatan} & : & Magang MBKM sebagai Software Engineer Intern di OPPO Manufacturing Indonesia \\
    \textbf{Proyek Terkait} & : & 360 Peer-to-Peer Grading System dan PQPR Software Project \\
\end{tabular}
\par\vspace{0.5\baselineskip}
\endgroup

Mata kuliah ini berfokus pada kemampuan melakukan pengujian terhadap unit dan modul perangkat lunak agar sistem yang dikembangkan memiliki perilaku yang sesuai dengan kebutuhan. Dalam kegiatan magang, pengujian dilakukan pada fitur-fitur yang berkaitan dengan proses evaluasi karyawan pada 360 Peer-to-Peer Grading System dan modul pendukung changeover produksi pada PQPR Software Project.

\section{Silabus}
Silabus mata kuliah mencakup pemahaman tujuan pengujian perangkat lunak, penyusunan skenario uji, pemeriksaan unit dan modul, validasi masukan dan keluaran, pelacakan kesalahan, debugging, serta dokumentasi hasil pengujian. Pengujian dilakukan untuk memastikan setiap bagian sistem bekerja sesuai rancangan sebelum digunakan dalam alur kerja yang lebih besar.

Dalam konteks proyek magang, pengujian tidak hanya dilakukan untuk menemukan kesalahan teknis, tetapi juga untuk memastikan bahwa fitur dapat digunakan dengan aman dan konsisten oleh pengguna internal perusahaan. Hal ini penting karena kesalahan pada modul evaluasi maupun modul produksi dapat memengaruhi keakuratan data dan kelancaran proses operasional.

\section{Aktivitas}
Pada proyek 360 Peer-to-Peer Grading System, aktivitas pengujian dilakukan terhadap alur pengisian penilaian, validasi data evaluasi, penyimpanan hasil, serta tampilan rekapitulasi. Skenario uji disusun berdasarkan kemungkinan penggunaan sistem, seperti pengguna mengisi data secara lengkap, pengguna melewati isian wajib, dan sistem menampilkan kembali hasil penilaian. Pengujian tersebut membantu memastikan bahwa proses evaluasi dapat berjalan terstruktur dan tidak menghasilkan data yang tidak valid.

Pada PQPR Software Project, pengujian dilakukan terhadap modul yang berkaitan dengan informasi produk, konfigurasi lini produksi, dan data kemampuan operator. Aktivitas pengujian mencakup pemeriksaan kesesuaian keluaran modul terhadap data masukan, pengecekan perubahan data, serta validasi bahwa informasi yang ditampilkan tetap konsisten setelah dilakukan pembaruan. Karena proyek ini berhubungan dengan proses manufaktur, pengujian perlu memperhatikan akurasi data dan keterlacakan perubahan.

Aktivitas lain yang dilakukan adalah debugging terhadap perilaku sistem yang tidak sesuai, melakukan pengecekan ulang setelah perbaikan, dan mendiskusikan hasil temuan dengan pembimbing maupun anggota tim. Setiap temuan pengujian digunakan sebagai dasar untuk memperbaiki fitur atau menyesuaikan logika modul agar lebih sesuai dengan kebutuhan pengguna.

\section{Konsep Dasar yang Berkaitan dengan Silabus}
Konsep dasar yang digunakan dalam pengujian unit adalah memastikan bagian terkecil dari sistem bekerja sesuai tanggung jawabnya. Unit yang diuji perlu memiliki masukan yang jelas, keluaran yang dapat diprediksi, dan penanganan kondisi tidak valid. Konsep ini diterapkan pada validasi form, pengolahan data, dan fungsi pendukung yang digunakan dalam fitur penilaian maupun modul produksi.

Pengujian modul berfokus pada interaksi beberapa unit dalam satu alur kerja. Modul yang tampak benar secara terpisah masih dapat menghasilkan masalah ketika dihubungkan dengan modul lain. Oleh karena itu, pengujian modul perlu memperhatikan alur data, perubahan status, hubungan antar fitur, dan kesesuaian keluaran dengan kebutuhan bisnis.

Konsep lain yang berkaitan adalah debugging dan regresi. Debugging dilakukan untuk menemukan penyebab perilaku yang tidak sesuai, sedangkan pengecekan regresi dilakukan untuk memastikan perbaikan tidak merusak fungsi yang sebelumnya sudah berjalan. Dalam proyek magang, kedua konsep tersebut penting karena fitur dikembangkan secara bertahap dan mengalami penyesuaian berdasarkan umpan balik.

\section{Hasil dan Pembahasan}
Hasil dari pelaksanaan mata kuliah ini adalah meningkatnya kemampuan menyusun dan menjalankan pengujian yang relevan terhadap modul perangkat lunak. Pada 360 Peer-to-Peer Grading System, pengujian membantu memastikan bahwa alur evaluasi karyawan berjalan sesuai rancangan, data yang wajib diisi dapat divalidasi, dan hasil penilaian dapat ditampilkan dengan konsisten. Hal ini mendukung keandalan sistem sebagai sarana pengumpulan umpan balik internal.

Pada PQPR Software Project, pengujian membantu memastikan bahwa modul yang digunakan untuk mendukung proses changeover produksi menampilkan informasi yang tepat dan merespons perubahan data dengan benar. Validasi terhadap data produk dan kemampuan operator menjadi bagian penting karena informasi tersebut digunakan untuk mendukung keputusan operasional. Pengujian yang dilakukan membantu mengurangi risiko kesalahan data pada sistem.

Pembahasan kegiatan menunjukkan bahwa pengujian perangkat lunak perlu dilakukan sejak proses pengembangan berlangsung, bukan hanya pada akhir pengerjaan. Dengan pengujian bertahap, kesalahan dapat ditemukan lebih awal dan perbaikan dapat dilakukan sebelum masalah menyebar ke modul lain. Pengujian juga menjadi sarana komunikasi teknis karena hasil temuan perlu dijelaskan kepada anggota tim agar perbaikan dapat dilakukan secara tepat.

\section{Kesimpulan}
Mata kuliah Internship: Pengujian Unit dan Modul Proyek terlaksana melalui aktivitas penyusunan skenario uji, validasi fitur, debugging, dan pengecekan ulang pada 360 Peer-to-Peer Grading System serta PQPR Software Project. Kegiatan ini memperkuat pemahaman bahwa kualitas perangkat lunak bergantung pada kemampuan menguji perilaku sistem secara sistematis. Pengujian yang baik membantu memastikan modul bekerja sesuai kebutuhan, data tetap konsisten, dan sistem lebih siap digunakan dalam lingkungan kerja perusahaan.
